\documentclass[11pt]{article}
\usepackage[margin=1in]{geometry}
\usepackage{amsmath,amssymb,amsthm}
\usepackage{enumitem}
\usepackage[hidelinks]{hyperref}

\title{Finite-\(p\) Hausdorff Spaceability for Prescribed Graph Dimension}
\author{Automatic ArXiv Shorter / Codex}
\date{12 June 2026}

\newtheorem*{theorem}{Theorem}
\newtheorem*{lemma}{Lemma}
\newtheorem*{proposition}{Proposition}
\newtheorem*{definition}{Definition}
\newtheorem*{remark}{Remark}

\newcommand{\cE}{\mathcal E}
\newcommand{\cH}{\mathcal H}
\newcommand{\NN}{\mathbb N}
\newcommand{\RR}{\mathbb R}
\newcommand{\cc}{\mathfrak c}
\newcommand{\dimH}{\dim_{\!H}}
\newcommand{\supp}{\operatorname{supp}}
\newcommand{\spanop}{\operatorname{span}}

\begin{document}
\maketitle

\section*{Problem}

In Liu--Shi--Zhang, \emph{On strong spaceability of continuous functions and fractal dimensions}, the following question appears at
\[
\texttt{downloads/extracted/2605.25037v2/Onspaceability.tex:1358}.
\]
For \(1<s\leq2\), put
\[
H_s[0,1]=\{f\in C[0,1]:\dimH G_f([0,1])=s\}.
\]
The question is whether \(H_s[0,1]\) is \((p,\cc)\)-spaceable for every finite integer \(3\leq p<\aleph_0\).

\begin{definition}
Let \(X\) be a topological vector space, \(A\subset X\), and \(p\in\NN\).  We say that \(A\) is \((p,\cc)\)-spaceable if, whenever \(E\) is a \(p\)-dimensional linear subspace of \(X\) such that
\[
E\subset A\cup\{0\},
\]
there is a closed linear subspace \(Y\subset X\) with
\[
E\subset Y\subset A\cup\{0\}
\]
and \(\dim Y=\cc\).
\end{definition}

\section*{Imported result from Liu--Shi--Zhang}

We use only the following special case of Liu--Shi--Zhang Proposition 3.4, taking \(K=[0,1]\) and \(m=1\).

\begin{proposition}[Liu--Shi--Zhang, Proposition 3.4 with \(m=1\)]
Let \(1<s\leq2\), let \(p\in\NN\), and let \(f_1,\ldots,f_p\in H_s[0,1]\) be linearly independent.  Assume
\[
\spanop\{f_1,\ldots,f_p\}\subset H_s[0,1]\cup\{0\}.
\]
Then there is a compact set \(K_0\subset[0,1]\) such that
\[
\dimH K_0=s-1
\]
and, for every nonzero \((a_1,\ldots,a_p)\in\RR^p\),
\[
\left.\sum_{i=1}^p a_i f_i\right|_{[0,1]\setminus K_0}
\in H_s([0,1]\setminus K_0).
\]
\end{proposition}

\section*{Auxiliary facts}

We shall use the following standard facts without further comment.

\begin{enumerate}[label=(F\arabic*),leftmargin=*]
\item Hausdorff dimension is countably stable:
\[
\dimH\left(\bigcup_{n=1}^{\infty} A_n\right)=\sup_n \dimH A_n.
\]
\item If \(K\subset\RR\) and \(f\in C(K)\), then
\[
\dimH G_f(K)\leq \dimH K+1.
\]
Indeed \(G_f(K)\subset K\times f(K)\), and \(f(K)\subset\RR\).
\item If \(K\subset[0,1]\) is compact and \(d=\dimH K\), then \(K\) has a full-dimensional point: there is \(x_0\in K\) such that
\[
\dimH(K\cap(x_0-r,x_0+r))=d,\qquad r>0.
\]
Otherwise finitely many neighborhoods of dimension \(<d\) would cover \(K\), contradicting countable stability.
\item If \(E\subset\RR\) is an uncountable compact set, then there exists \(u\in C(E)\) such that
\[
\dimH G_u(E)=\dimH E+1.
\]
This is the prevalent maximal-graph-dimension theorem cited as Lemma 2.4 in Liu--Shi--Zhang.
\end{enumerate}

\section*{The compact \(c_0\)-subspace lemma}

\begin{lemma}\label{lem:c0-compact}
Let \(K\subset[0,1]\) be compact and let
\[
d=\dimH K>0.
\]
Then there is a closed subspace \(X_K\subset C[0,1]\), isomorphic to \(c_0\), with the following properties.
\begin{enumerate}[label=(\roman*),leftmargin=*]
\item For every nonzero \(h\in X_K\),
\[
\dimH G_h(K)=d+1.
\]
\item Every \(h\in X_K\) is affine on each connected component of \([0,1]\setminus K\).
\end{enumerate}
\end{lemma}

\begin{proof}
Choose a full-dimensional point \(x_0\in K\).  We first construct disjoint pieces of \(K\) whose dimensions tend to \(d\).

Let \((r_m)\) be a strictly decreasing sequence of positive numbers tending to \(0\), chosen so fast that the annuli
\[
A_m=K\cap\bigl((x_0-r_m,x_0+r_m)\setminus [x_0-r_{m+1},x_0+r_{m+1}]\bigr)
\]
are pairwise disjoint.  Since
\[
K\cap(x_0-r_N,x_0+r_N)
\subset \{x_0\}\cup\bigcup_{m\geq N} A_m,
\]
and the left-hand side has Hausdorff dimension \(d\), countable stability gives
\[
\sup_{m\geq N}\dimH A_m=d
\]
for every \(N\).

Therefore, after passing to a subsequence of annuli and then using inner regularity of Hausdorff dimension, we may choose compact sets \(E_n\subset K\) and open intervals \(I_n\) such that
\[
E_n\subset I_n,\qquad \overline{I_n}\cap\overline{I_m}=\emptyset\ (n\neq m),\qquad I_n\to x_0,
\]
and
\[
d_n:=\dimH E_n>d-\frac1n.
\]
Discard finitely many terms if necessary, so that \(d_n>0\) for every \(n\).  Then each \(E_n\) is uncountable, and \(d_n\to d\).

By (F4), for each \(n\) there is \(\varphi_n\in C(E_n)\) such that
\[
\dimH G_{\varphi_n}(E_n)=d_n+1.
\]
Multiplying by a nonzero scalar if needed, assume \(\|\varphi_n\|_\infty=1\).

Since \(E_n\) and \(K\setminus I_n\) are disjoint closed subsets of \(K\), Tietze's extension theorem gives a function \(\psi_n\in C(K)\) such that
\[
\psi_n|_{E_n}=\varphi_n,\qquad
\psi_n|_{K\setminus I_n}=0,\qquad
\|\psi_n\|_\infty=1.
\]
The supports of the \(\psi_n\)'s are pairwise disjoint subsets of \(K\), because the intervals \(I_n\) are pairwise disjoint.

Partition \(\NN\) into pairwise disjoint infinite sets
\[
\NN=B_1\cup B_2\cup\cdots.
\]
For each \(j\), choose coefficients \(c_{j,n}\in(0,1]\), \(n\in B_j\), such that
\[
\sup_{n\in B_j}c_{j,n}=1
\quad\hbox{and}\quad
c_{j,n}\to0\quad(n\to\infty,\ n\in B_j).
\]
For example, take one coefficient equal to \(1\), and make all later coefficients tend to \(0\).

Define \(v_j\in C(K)\) by
\[
v_j(x)=
\begin{cases}
c_{j,n}\psi_n(x),& x\in K\cap I_n\hbox{ for some }n\in B_j,\\
0,&\hbox{otherwise}.
\end{cases}
\]
The only possible accumulation point of the intervals \(I_n\) is \(x_0\), and the coefficients \(c_{j,n}\) tend to \(0\) along \(B_j\).  Hence each \(v_j\) is continuous at \(x_0\), and it is plainly continuous elsewhere.  Also \(\|v_j\|_\infty=1\).

Now define
\[
T:c_0\longrightarrow C(K),\qquad
T((a_j)_{j\geq1})=\sum_{j=1}^\infty a_jv_j.
\]
This is a pointwise locally finite sum away from \(x_0\), and continuity at \(x_0\) follows from \(a_j\to0\) together with \(v_j(x_0)=0\).  Since the supports of the \(v_j\)'s are pairwise disjoint,
\[
\|T(a)\|_\infty=\sup_j |a_j|.
\]
Thus \(T\) is an isometric embedding and \(T(c_0)\) is a closed subspace of \(C(K)\).

If \(a=(a_j)\neq0\), choose \(j_0\) with \(a_{j_0}\neq0\).  For every \(n\in B_{j_0}\), on \(E_n\) we have
\[
T(a)=a_{j_0}c_{j_0,n}\varphi_n.
\]
Nonzero scalar multiplication does not change Hausdorff dimension, so
\[
\dimH G_{T(a)}(E_n)=d_n+1,\qquad n\in B_{j_0}.
\]
Therefore
\[
\dimH G_{T(a)}(K)\geq
\sup_{n\in B_{j_0}}(d_n+1)=d+1.
\]
The reverse inequality is (F2), so
\[
\dimH G_{T(a)}(K)=d+1.
\]

Finally, let \(\cE_K:C(K)\to C[0,1]\) be the canonical linear extension obtained by affine interpolation on every connected component of \([0,1]\setminus K\).  This operator is an isometric embedding.  Put
\[
X_K=\cE_K(T(c_0)).
\]
Then \(X_K\) is closed in \(C[0,1]\), isomorphic to \(c_0\), satisfies the same graph-dimension conclusion on \(K\), and is affine on every component of \([0,1]\setminus K\).
\end{proof}

\section*{Main theorem}

\begin{theorem}
Let \(1<s\leq2\).  Then \(H_s[0,1]\) is \((p,\cc)\)-spaceable for every finite \(p\in\NN\).
\end{theorem}

\begin{proof}
Fix \(p\in\NN\).  Let \(E\subset C[0,1]\) be a \(p\)-dimensional subspace such that
\[
E\subset H_s[0,1]\cup\{0\}.
\]
Choose a basis \(f_1,\ldots,f_p\) of \(E\).  Applying the imported Liu--Shi--Zhang proposition gives a compact set \(K_0\subset[0,1]\) such that
\[
\dimH K_0=s-1
\]
and
\[
w|_{[0,1]\setminus K_0}\in H_s([0,1]\setminus K_0),
\qquad 0\neq w\in E.
\]
Set \(d=s-1\).  Since \(s>1\), we have \(d>0\).  Apply Lemma \ref{lem:c0-compact} to \(K_0\), and let \(X_0=X_{K_0}\).  Thus \(X_0\) is closed, \(\dim X_0=\cc\), every nonzero \(h\in X_0\) has
\[
\dimH G_h(K_0)=s,
\]
and every \(h\in X_0\) is affine on the connected components of \([0,1]\setminus K_0\).

We first prove that
\[
E\cap X_0=\{0\}.
\]
Suppose \(0\neq w\in E\cap X_0\).  Since \(w\in E\), the choice of \(K_0\) gives
\[
\dimH G_w([0,1]\setminus K_0)=s.
\]
Since \(w\in X_0\), the function \(w\) is affine on each connected component of \([0,1]\setminus K_0\).  The set \([0,1]\setminus K_0\) is a countable union of open intervals, and the graph of an affine function on each interval has Hausdorff dimension \(1\).  Hence
\[
\dimH G_w([0,1]\setminus K_0)\leq1,
\]
contradicting \(s>1\).  Thus the intersection is trivial.

Let
\[
Y=E\oplus X_0.
\]
Because \(E\) is finite-dimensional and \(X_0\) is closed, the direct sum \(Y\) is a closed subspace of \(C[0,1]\).  Moreover \(E\subset Y\) and \(\dim Y=\cc\).

It remains to show that
\[
Y\subset H_s[0,1]\cup\{0\}.
\]
Let \(0\neq y\in Y\).  Write uniquely
\[
y=w+h,\qquad w\in E,\ h\in X_0.
\]

Assume first that \(w\neq0\).  Let
\[
U=[0,1]\setminus K_0.
\]
On every connected component of \(U\), the function \(h\) is affine.  Therefore, on each such component \(J\), the map
\[
(x,t)\mapsto (x,t+h(x))
\]
is bi-Lipschitz on \(J\times\RR\).  Hence
\[
\dimH G_{w+h}(J)=\dimH G_w(J).
\]
Using countable stability over the components \(J\) of \(U\), we get
\[
\dimH G_y(U)=\dimH G_w(U)=s.
\]
On the compact set \(K_0\), the general upper estimate gives
\[
\dimH G_y(K_0)\leq \dimH K_0+1=s.
\]
Thus
\[
\dimH G_y([0,1])
=\max\{\dimH G_y(K_0),\dimH G_y(U)\}=s.
\]

Assume now that \(w=0\).  Then \(y=h\neq0\), and Lemma \ref{lem:c0-compact} gives
\[
\dimH G_y(K_0)=s.
\]
The same upper estimate gives \(\dimH G_y([0,1])\leq s\), while the previous display gives the reverse inequality.  Hence again
\[
\dimH G_y([0,1])=s.
\]

So every nonzero element of \(Y\) belongs to \(H_s[0,1]\).  Since the initial \(p\)-dimensional subspace \(E\) was arbitrary, \(H_s[0,1]\) is \((p,\cc)\)-spaceable.
\end{proof}

\section*{Consequence}

Liu--Shi--Zhang also prove that \(H_s[0,1]\) is not \((\alpha,\beta)\)-spaceable for any \(\alpha\geq\aleph_0\), regardless of \(\beta\).  Therefore the theorem above gives the positive side for every possible finite initial dimension \(p\), and the negative side for every infinite initial dimension was already known in their paper.

\section*{Human-check checklist}

\begin{itemize}[leftmargin=*]
\item Verify the extraction of the \(m=1\) case of Liu--Shi--Zhang Proposition 3.4.
\item Verify the annulus selection in Lemma \ref{lem:c0-compact}; it is the only place where dimension concentration at one point matters.
\item Verify the use of affine interpolation \(\cE_K\), especially that it is isometric and affine on complementary intervals.
\item Verify that the countable component argument on \([0,1]\setminus K_0\) replaces the ``nearly locally Lipschitz'' lemma without requiring compactness of the domain.
\end{itemize}

\end{document}
